% BHU PYQ -- Manually transcribed & error-corrected
% Paper: BCH-111 : Financial Accounting
% Exam: B.Com. (Hons.) Semester I Examination, 2016-17

\documentclass[11pt,a4paper]{article}
\usepackage[a4paper,top=2.5cm,bottom=2.5cm,left=2.8cm,right=2.8cm,headheight=14pt]{geometry}
\usepackage{amsmath,amssymb}
\usepackage[T1]{fontenc}
\usepackage[utf8]{inputenc}
\usepackage{lmodern,microtype,fancyhdr,enumitem,hyperref,lastpage,parskip}

\hypersetup{
    colorlinks=true,
    urlcolor=black,
    linkcolor=black,
    pdftitle={BCH-111: Financial Accounting -- BHU B.Com. (Hons.) Sem I 2016-17}
}

\pagestyle{fancy}
\fancyhf{}
\renewcommand{\headrulewidth}{0.4pt}
\renewcommand{\footrulewidth}{0.4pt}
\fancyhead[L]{\small\textit{Banaras Hindu University}}
\fancyhead[R]{\small\textit{BCH-111: Financial Accounting}}
\fancyfoot[C]{\small Page \thepage\ of \pageref{LastPage}}

\newlist{parts}{enumerate}{2}
\setlist[parts,1]{label=(\alph*),leftmargin=*,itemsep=5pt,topsep=3pt}
\setlist[parts,2]{label=(\roman*),leftmargin=*,itemsep=3pt,topsep=2pt}
\newcommand{\pts}[1]{\hfill\textit{[#1]}}

\begin{document}

\begin{center}
  {\large\bfseries BANARAS HINDU UNIVERSITY}\\[2pt]
  {\normalsize B.Com. (Hons.) --- Semester I Examination, 2016-17}\\[6pt]
  \rule{0.8\linewidth}{0.5pt}\\[6pt]
  {\large\bfseries Paper No. BCH-111: Financial Accounting}\\[4pt]
  {\normalsize Time: 3 Hours\hfill Maximum Marks: 70}
\end{center}

\rule{\linewidth}{0.4pt}

\smallskip
\noindent\textit{Instructions: Answer all questions. All questions carry equal marks (14 marks each).}

\bigskip

\noindent\textbf{Question 1.}
What do you understand by Capital and Revenue Expenditure? Why is it essential to distinguish between the two? Explain any six revenue expenses which become Capital Expenditure.\pts{14}

\centerline{\textbf{OR}}

Journalize the following transactions of Pankaj Prakash Chandra:\pts{14}
\begin{parts}
  \item Bought goods worth Rs. 1,000 from Usha and supplied them to Kiran at Rs. 1,500.
  \item An employee stole away: Cash Rs. 1,500 and Goods Rs. 500.
  \item Goods valuing Rs. 15,000 insured for Rs. 10,000 @ 5\% as premium paid.
  \item Purchased goods from Samuel Rs. 50,000 less trade discount @ 20\%, plus VAT @ 10\%.
  \item Sold goods costing Rs. 6,000 to Mohan issued at 10\% above cost less 5\% trade discount.
  \item Paid Rs. 2,500 as wages on installation of new Machine.
  \item Sold to Arvind at a trade discount @ 5\% and cash discount of 1\% goods of the list price of Rs. 20,000. He paid 60\% in cash.
\end{parts}

\medskip\rule{\linewidth}{0.2pt}

\noindent\textbf{Question 2.}
Explain how one-sided and double-sided errors are corrected? Give suitable examples also.\pts{14}

\centerline{\textbf{OR}}

A manufacturing firm purchased certain machinery on 1st January, 2012 for Rs. 19,400 and spent Rs. 600 on its installation. On 1st July in the same year, additional machinery costing Rs. 10,000 was purchased. On July 1st 2014, the machinery purchased on 1st January, 2012 having become obsolete, was auctioned for Rs. 8,000 and on the same date fresh machinery was purchased at a cost of Rs. 15,000.

Depreciation was provided annually on 31st December at the rate of 10\% per annum on the original cost of the assets. In 2015, however, the firm changed this method to one of writing off 15\% p.a. on the written down value.

Give machinery account as it would stand at the end of each year from 2012 to 2016.\pts{14}

\medskip\rule{\linewidth}{0.2pt}

\noindent\textbf{Question 3.}
\begin{parts}
  \item Explain the Fixed and Fluctuating Capital Accounts.\pts{7}
  \item What is hidden goodwill? How is it computed?\pts{7}
\end{parts}

\centerline{\textbf{OR}}

\begin{parts}
  \item A and B are partners sharing profits and losses equally. They admit C for $1/4^{\text{th}}$ share. C pays only Rs. 5,000 for premium out of his premium of Rs. 9,000. Goodwill already appears in the balance sheet at Rs. 30,000. Give journal entries.\pts{7}
  \item X and Y have been in partnership sharing profits in the ratio of 7:3. Z is admitted as a partner. X surrenders $1/7^{\text{th}}$ of his share and Y surrenders $1/3^{\text{rd}}$ of his share in favour of Z. Calculate the new profit sharing ratio and sacrificing ratio.\pts{7}
\end{parts}

\medskip\rule{\linewidth}{0.2pt}

\noindent\textbf{Question 4.}
Explain the different methods of recording a Joint Life Policy in the books of accounts regarding the retirement of a partner.\pts{14}

\centerline{\textbf{OR}}

X, Y and Z are partners in a business sharing profits and losses in the ratio of 3:2:1. On 31st December, 2013, Z died. His capital account after all necessary adjustments shows a balance of Rs. 1,09,500. It is agreed to transfer this amount to his executor account and then immediately he should be paid Rs. 29,500 in cash, Rs. 20,000 through bills payable and the balance amount in three equal yearly installments with interest @ 8\% p.a. You are required to show Z's Executor account till it is completely paid up.\pts{14}

\medskip\rule{\linewidth}{0.2pt}

\noindent\textbf{Question 5.}
What is a Realization Account? Why is this account prepared? Give a specimen of Realization Account.\pts{14}

\centerline{\textbf{OR}}

Jai, Shiv and Shanker entered into business as a partnership on April 1, 2015 sharing profits \& losses in the proportion of 4:3:2 respectively and with capitals of Rs. 30,000, Rs. 20,000 and Rs. 10,000. Their assets and liabilities on March 31, 2016, the date on which they decided to wind-up their affairs, were as follows: office fixtures Rs. 1,000, debtors Rs. 28,000, bills receivable Rs. 5,000, stock Rs. 45,000, sundry creditors Rs. 30,000 and bills payable Rs. 4,000.

Jai agreed to take over the stock at a discount of 10\% and pay off the bills payable. Shiv agreed to take over the book debts at a discount of 20\% and pay off the creditors. Shankar took over the bills receivable at Rs. 4,877 and office fixtures at a depreciation of 10\%. 5\% interest is credited to each partner on his capital.

Prepare the Realization Account and Partners' Capital Accounts.\pts{14}

\vfill
\rule{\linewidth}{0.4pt}
\begin{center}\small
\href{https://bhu-exam.vercel.app/}{Click here to visit our website}\\[4pt]
\href{https://me-aryan.vercel.app/}{Click here to visit our contributor}\\[4pt]
\href{https://quashan-me.vercel.app/}{Click here to visit our contributor}
\end{center}

\end{document}
